\section{The two-letter criterion}\label{sec:classification}

Fix an irrational \(\alpha\in(0,1)\) and a window length
\(\beta\in(0,1)\).  On the circle \(\TT=\mathbb{R}/\mathbb{Z}\), define
\[
s_j(x):=\mathbf{1}_{[0,\beta)}(x+j\alpha)
\qquad(j\geq0).
\]
The realized length-\(m\) language is
\[
S_m(\alpha,\beta):=
\left\{s_0(x)\cdots s_{m-1}(x):
x\in\TT\text{ on a set of positive Lebesgue measure}\right\}.
\]
Equivalently, \(S_m(\alpha,\beta)\) is the support of the block law obtained
by choosing \(x\) according to Lebesgue measure.  Requiring positive measure
removes the irrelevant choices caused by endpoints of the half-open window.

We begin with the circle observation that will be used twice.

\begin{lemma}[Equal-length interval overlap]\label{lem:interval-overlap}
Let \(I\) and \(J\) be half-open circle intervals of the same length
\(\ell\in(0,1)\), and suppose their initial points differ by
\(\alpha\) modulo \(1\).  Then
\[
\leb(I\cap J)>0
\quad\Longleftrightarrow\quad
\ell>\delta(\alpha),
\qquad
\delta(\alpha)=\min\{\alpha,1-\alpha\}.
\]
\end{lemma}

\begin{proof}
Choose the shorter of the two orientations between the initial points.
Its length is \(\delta(\alpha)\in(0,1/2)\).  In that orientation the two
equal half-open arcs are disjoint up to endpoints exactly when their common
length is at most the displacement.  Endpoints have Lebesgue measure zero,
so the intersection has positive measure exactly when
\(\ell>\delta(\alpha)\).
\end{proof}

\begin{theorem}[Sharp classification of collision-free windows]
\label{thm:window-injectivity-classification}
Let \(\alpha\in(0,1)\) be irrational, let \(\beta\in(0,1)\), and put
\(\delta=\delta(\alpha)\).  The following are equivalent.
\begin{enumerate}
\item For every \(m\geq1\), the restriction
\[
\Fold_m|_{S_m(\alpha,\beta)}:S_m(\alpha,\beta)\longrightarrow X_m
\]
is injective.
\item The restriction of \(\Fold_2\) to \(S_2(\alpha,\beta)\) is injective.
\item The window length satisfies
\[
\beta\in(0,\delta]\cup[1-\delta,1).
\]
\end{enumerate}
More precisely:
\begin{enumerate}
\item[(a)] If \(\beta\leq\delta\), then
\(S_m(\alpha,\beta)\subset X_m\), and \(\Fold_m\) is the identity on the
realized language.
\item[(b)] If \(\beta\geq1-\delta\), then the coordinatewise complement of
every word in \(S_m(\alpha,\beta)\) belongs to \(X_m\), and \(\Fold_m\) is
injective on the realized language.
\item[(c)] If \(\delta<\beta<1-\delta\), then
\(00,11\in S_2(\alpha,\beta)\) and
\[
\Fold_2(00)=\Fold_2(11)=00.
\]
\end{enumerate}
\end{theorem}

\begin{proof}
Put \(I=[0,\beta)\).  The two-letter word \(11\) occurs on
\[
I\cap(I-\alpha),
\]
where the two intervals have length \(\beta\) and their initial points
differ by \(\alpha\).  Lemma~\ref{lem:interval-overlap} gives
\[
11\in S_2(\alpha,\beta)
\quad\Longleftrightarrow\quad
\beta>\delta.
\]
Likewise, \(00\) occurs on the intersection of two translates of the
complementary interval.  Those intervals have length \(1-\beta\), so
\[
00\in S_2(\alpha,\beta)
\quad\Longleftrightarrow\quad
1-\beta>\delta.
\]

Suppose first that \(\beta\leq\delta\).  Then \(11\) does not occur with
positive measure.  If a realized word of any length contained adjacent
ones, the cylinder on which it occurs would be contained in a translate of
the two-letter \(11\)-cylinder, a null set.  Hence
\[
S_m(\alpha,\beta)\subset X_m
\qquad(m\geq1).
\]
Proposition~\ref{prop:fold-residue-characterization} says that the fold
fixes \(X_m\), proving (a) and injectivity in the sparse regime.

Now suppose that \(\beta\geq1-\delta\).  Then \(00\) does not occur with
positive measure.  For \(u=u_1\cdots u_m\in S_m(\alpha,\beta)\), let
\[
\widetilde u:=(1-u_1)\cdots(1-u_m).
\]
The absence of \(00\) means that \(\widetilde u\in X_m\).  If \(u\neq v\),
then \(\widetilde u\neq\widetilde v\), and the bijection in
Proposition~\ref{prop:fold-residue-characterization} gives
\[
0<
\left|N_m(\widetilde u)-N_m(\widetilde v)\right|
<F_{m+2}.
\]
The sum \(N_m(u)+N_m(\widetilde u)=\sum_{j=1}^mF_{j+1}\) is independent of
\(u\).  Therefore
\[
N_m(u)-N_m(v)
=-\bigl(N_m(\widetilde u)-N_m(\widetilde v)\bigr).
\]
This nonzero difference has absolute value less than \(F_{m+2}\), so the
values \(N_m(u)\) and \(N_m(v)\) are distinct modulo \(F_{m+2}\).
The residue characterization now gives
\(\Fold_m(u)\neq\Fold_m(v)\), proving (b).

Finally, suppose that \(\delta<\beta<1-\delta\).  The two overlap tests show
that both \(11\) and \(00\) belong to \(S_2(\alpha,\beta)\).  Since
\[
N_2(00)=0,\qquad
N_2(11)=F_2+F_3=F_4,
\]
Proposition~\ref{prop:fold-residue-characterization} yields
\(\Fold_2(00)=\Fold_2(11)=00\).  Thus condition (2) fails.

We have proved that (3) implies (1), while (1) trivially implies (2), and
failure of (3) implies failure of (2).  This establishes the equivalence and
the three refinements.
\end{proof}

The customary Sturmian window lies exactly on one of the two boundary
points, depending on the orientation of the rotation.

\begin{corollary}[Sturmian blocks are folded without collision]
\label{cor:sturmian-no-collision}
Assume \(\beta=\alpha\).  For every \(m\geq1\), the restriction
\[
\Fold_m|_{S_m(\alpha,\alpha)}
\]
is injective.
\end{corollary}

\begin{proof}
If \(\alpha<1/2\), then
\(\beta=\alpha=\delta(\alpha)\).  If \(\alpha>1/2\), then
\(\beta=\alpha=1-\delta(\alpha)\).  Irrationality excludes
\(\alpha=1/2\), so Theorem~\ref{thm:window-injectivity-classification}
applies in either case.
\end{proof}

\begin{remark}[Higher-block presentation]\label{rem:higher-block}
Let \(S_m(\alpha,\alpha)\) be used as the alphabet of the ordinary
\(m\)-block presentation of the Sturmian shift.  By
Corollary~\ref{cor:sturmian-no-collision}, the restriction of \(\Fold_m\)
is a bijection from this realized alphabet onto its image.  Applying that
bijection coordinatewise merely relabels the ordinary higher-block
presentation.  It is therefore a one-block conjugacy onto the folded image,
with the coordinatewise inverse.  This is an immediate symbolic consequence
of injectivity, not a broader circle-dynamical theorem.
\end{remark}

\begin{example}[The three regimes]\label{ex:three-regimes}
Take the golden rotation
\[
\alpha=\frac{\sqrt5-1}{2}.
\]
Then
\[
\delta(\alpha)=1-\alpha=\frac{3-\sqrt5}{2}.
\]
Three choices of \(\beta\) display all possibilities.
\begin{enumerate}
\item If \(\beta=1/4<\delta(\alpha)\), then \(11\) has zero measure.
Every realized word already belongs to \(X_m\), so every fold is the
identity on the realized language.
\item If \(\beta=1/2\), then
\(\delta(\alpha)<\beta<1-\delta(\alpha)\).  Both \(00\) and \(11\) occur
with positive measure, and they collide under \(\Fold_2\).
\item If \(\beta=3/4>1-\delta(\alpha)\), then \(00\) has zero measure.
Complementing every realized word puts it in \(X_m\), and the residue
argument in the theorem makes every fold injective on the realized
language.
\end{enumerate}
These alternatives reflect exactly the interval-overlap trichotomy.
At the endpoint \(\beta=\alpha=1-\delta(\alpha)\), the coding is Sturmian
and belongs to the dense collision-free regime.
\end{example}

\begin{remark}[Scope]\label{rem:scope}
The theorem does not use discrepancy estimates, empirical histograms,
Parry measures, or a transfer to conjugate circle maps.  Its
all-resolution content is precisely the two-letter overlap test together
with Proposition~\ref{prop:fold-residue-characterization} in the dense
case.  The endpoint equalities are included because the language is defined
by positive-measure occurrence: at equality the relevant interval
intersection consists only of endpoints and contributes no word to the
support.
\end{remark}
